\chapter{Cardiovascular System}

\section*{Definition and Overview}
The cardiovascular system consists of the heart, the blood vessels, and the blood they carry. It forms a closed circuit through which oxygen, nutrients, hormones, and waste products are transported around the body, and it is central to almost every function of human physiology. The heart is a muscular pump divided into four chambers, and it beats roughly 100,000 times each day to propel blood through the lungs, where it picks up oxygen, and through the rest of the body, where oxygen and nutrients are delivered to tissues. Understanding the normal structure and function of the cardiovascular system is fundamental to clinical medicine, because disorders of the heart and blood vessels are the leading cause of death worldwide.

\section*{Epidemiology}
Cardiovascular disease is the leading cause of death globally, accounting for around one in three deaths. The most common manifestations are coronary artery disease, which causes heart attacks, and cerebrovascular disease, which causes stroke. Hypertension, or high blood pressure, affects more than one in four adults and is the single biggest contributor to cardiovascular risk. The burden increases with age and is higher in men, in lower-income settings, and in people with diabetes, obesity, and smoking histories. Many cardiovascular diseases are preventable, and declining smoking rates and better blood pressure and cholesterol control have substantially reduced death rates in many high-income countries over recent decades.

\section*{Aetiology and Pathophysiology}
The cardiovascular system can be affected by a wide range of disease processes, each disrupting normal anatomy and physiology:
\begin{itemize}
    \item \textbf{Atherosclerosis:} the progressive accumulation of fatty plaques within artery walls, narrowing vessels and predisposing to heart attacks and stroke
    \item \textbf{Hypertension:} chronically elevated blood pressure that damages artery walls and forces the heart to work harder
    \item \textbf{Heart failure:} the inability of the heart to pump enough blood for the body's needs, from muscle damage, chronic pressure overload, or stiffening
    \item \textbf{Arrhythmias:} disorders of the heart's electrical system producing fast, slow, or irregular rhythms
    \item \textbf{Valvular disease:} narrowing or leakage of the heart valves, most commonly aortic stenosis and mitral regurgitation
    \item \textbf{Congenital heart disease:} structural abnormalities present from birth
    \item \textbf{Inflammatory and infective conditions:} such as myocarditis, pericarditis, and infective endocarditis
    \item \textbf{Venous and thromboembolic disease:} deep vein thrombosis and pulmonary embolism
\end{itemize}

\section*{Clinical Features}
Cardiovascular disease produces symptoms that reflect the affected part of the system:
\begin{itemize}
    \item Chest pain, pressure, or tightness, often on exertion (angina) or at rest with a heart attack
    \item Shortness of breath with activity or when lying flat
    \item Swelling of the ankles, feet, or abdomen from fluid retention
    \item Palpitations, or a fluttering, racing, or irregular heartbeat
    \item Dizziness, light-headedness, or fainting
    \item Fatigue and reduced exercise tolerance
    \item Leg pain on walking that eases with rest, suggesting peripheral arterial disease
    \item Sudden weakness, speech difficulty, or facial droop, suggesting stroke
\end{itemize}

\section*{Differential Diagnosis}
Symptoms of cardiovascular disease overlap with many other conditions:
\begin{itemize}
    \item Chest pain: gastro-oesophageal reflux, musculoskeletal pain, anxiety, pneumonia, or pulmonary embolism
    \item Breathlessness: lung disease, anaemia, obesity, or anxiety
    \item Palpitations: caffeine, alcohol, anxiety, thyroid disease, or anaemia
    \item Leg swelling: kidney disease, liver disease, venous insufficiency, or medication effects
    \item Fainting: low blood sugar, dehydration, neurological conditions, or low blood pressure on standing
\end{itemize}

\section*{When to Seek Medical Care}
Anyone with new, severe, or worsening chest pain, breathlessness, palpitations, or swelling should see a doctor promptly. Risk-factor checks for blood pressure and cholesterol are recommended in adulthood, particularly from age 45 and earlier in people with a strong family history or other risk factors.
\textbf{Call 000 (or local emergency services) for:}
\begin{itemize}
    \item Chest pain or pressure lasting more than a few minutes, or at rest
    \item Pain spreading to the arm, jaw, neck, or back
    \item Sudden severe breathlessness, sweating, or feeling faint
    \item Signs of stroke: facial droop, arm weakness, or speech difficulty
    \item Collapse or loss of consciousness
\end{itemize}

\section*{Diagnosis and Investigations}
The cardiovascular system is assessed through a combination of clinical and technological methods:
\begin{itemize}
    \item \textbf{History and examination:} symptoms, risk factors, family history, and examination of the pulse, blood pressure, heart, lungs, and leg vessels
    \item \textbf{Blood pressure measurement:} the basic screening test for hypertension
    \item \textbf{Blood tests:} cholesterol, blood glucose, kidney function, and cardiac markers such as troponin when a heart attack is suspected
    \item \textbf{Electrocardiogram (ECG):} records the heart's electrical activity to detect heart attacks, arrhythmias, and strain patterns
    \item \textbf{Echocardiogram:} ultrasound imaging of heart structure, valve function, and pumping strength
    \item \textbf{Exercise stress testing:} assesses the heart's response to exertion and provokes symptoms
    \item \textbf{Ambulatory monitoring:} portable ECG devices worn for days to capture intermittent arrhythmias
    \item \textbf{Imaging:} coronary angiography, CT, MRI, or nuclear studies for detailed assessment of vessels and muscle
\end{itemize}

\section*{Management}
Management of cardiovascular disease ranges from prevention to advanced treatment:
\begin{itemize}
    \item \textbf{Healthy lifestyle:} regular physical activity, a heart-healthy diet, weight control, and stopping smoking
    \item \textbf{Medicines:} drugs to lower blood pressure and cholesterol, thin the blood, control heart rate and rhythm, and relieve fluid overload
    \item \textbf{Procedures:} angioplasty and stenting to open narrowed arteries
    \item \textbf{Surgery:} bypass grafting, valve repair or replacement, and surgery for structural problems
    \item \textbf{Devices:} pacemakers and implantable defibrillators for rhythm disorders
    \item \textbf{Cardiac rehabilitation:} structured exercise and education after cardiac events
    \item \textbf{Acute care:} rapid treatment of heart attacks and stroke in hospital settings
\end{itemize}

\section*{Complications}
Untreated or poorly controlled cardiovascular disease can lead to serious consequences:
\begin{itemize}
    \item Heart attack with permanent damage to heart muscle
    \item Heart failure with breathlessness, swelling, and repeated hospital admissions
    \item Stroke causing lasting disability or death
    \item Sudden cardiac death from dangerous arrhythmias
    \item Peripheral arterial disease with leg pain, poor healing, or amputation
    \item Kidney damage from hypertension and reduced blood flow
    \item Blood clots in the legs or lungs
    \item Reduced quality of life, depression, and cognitive decline
\end{itemize}

\section*{Prognosis and Outcomes}
Outcomes depend on the specific condition, its severity, and how early it is detected and treated. Many cardiovascular diseases can be controlled for decades with modern medication, and survival after heart attack has improved dramatically in recent decades because of rapid treatment and secondary prevention. Prognosis is generally better when risk factors are controlled, medicines are taken reliably, and follow-up is consistent. With a heart-healthy lifestyle, much of the disease burden is preventable in the first place.

\section*{Prevention and Self-Care}
\begin{itemize}
    \item Stay physically active, aiming for at least 150 minutes of moderate activity each week
    \item Eat a diet rich in vegetables, fruit, whole grains, fish, and healthy fats, and low in salt and saturated fat
    \item Maintain a healthy weight and avoid excess alcohol
    \item Do not smoke, and avoid exposure to second-hand smoke
    \item Have blood pressure, cholesterol, and blood glucose checked regularly
    \item Manage diabetes and other chronic conditions under medical supervision
    \item Know the warning signs of heart attack and stroke and act fast
\end{itemize}

\section*{Sources and Further Reading}
The following sources provide authoritative, up-to-date information on the cardiovascular system and cardiovascular disease.
\begin{itemize}
    \item World Health Organization. \textit{Cardiovascular diseases fact sheet.} https://www.who.int/
    \item NHS. \textit{Cardiovascular disease overview.} https://www.nhs.uk/conditions/cardiovascular-disease/
    \item Heart Foundation Australia. \textit{Heart disease information and prevention.} https://www.heartfoundation.org.au/
    \item American Heart Association. \textit{Heart and cardiovascular disease information.} https://www.heart.org/
    \item National Heart, Lung, and Blood Institute. \textit{Cardiovascular health information.} https://www.nhlbi.nih.gov/health
    \item MedlinePlus. \textit{Heart and vascular health.} https://medlineplus.gov/heartandvascular.html
\end{itemize}
